\documentclass[12pt,a4paper]{amsart}

\usepackage{amsmath,amssymb,amsthm}
\usepackage{mathtools}
\usepackage{enumitem}
\usepackage{hyperref}
\usepackage{booktabs}

%%%─── Theorem environments
\newtheorem{theorem}{Theorem}[section]
\newtheorem{lemma}[theorem]{Lemma}
\newtheorem{proposition}[theorem]{Proposition}
\newtheorem{corollary}[theorem]{Corollary}
\newtheorem{conjecture}[theorem]{Conjecture}
\newtheorem{hypothesis}[theorem]{Hypothesis}
\theoremstyle{definition}
\newtheorem{definition}[theorem]{Definition}
\newtheorem{problem}[theorem]{Open Problem}
\theoremstyle{remark}
\newtheorem{remark}[theorem]{Remark}
\newtheorem*{remark*}{Remark}

%%%─── Macros
\newcommand{\Phihat}{\widehat{\Phi}}
\newcommand{\norm}[1]{\|#1\|}
\newcommand{\ip}[2]{\langle #1,\, #2\rangle}
\newcommand{\Hc}{\mathcal{H}_c}
\newcommand{\Hlim}{\mathcal{H}_{\lim}}
\newcommand{\HSOT}{\mathcal{H}^{\mathrm{SOT}}}
\newcommand{\Hspectral}{H_{\lim}^{(\mathrm{sp})}}
\newcommand{\Dom}{\mathrm{Dom}}
\newcommand{\sigapp}{\sigma_{\mathrm{app}}}

\subjclass[2020]{47A10, 47B25, 42C05, 11M06}
\keywords{Prolate spheroidal wave functions, limiting PSWF operator,
spectral inclusion, semiclassical transfer, spectral density,
Riemann zeta function}

\begin{document}

\title[Spectral Structure of the Limiting PSWF Operator]{%
  Spectral Structure and Density\\
  for the Limiting PSWF Operator}

\author{[Author]}
\address{[Address]}
\email{[Email]}

\maketitle

\begin{abstract}
We continue the operator-theoretic analysis of the limiting PSWF
concentration operator $\Hlim$, established in Papers~IX--X.
The operator $\Hlim$ is the SOT-limit of $\Hc$ (Paper~X);
its eigenstructure is an orthonormal system
$\{\Phi_n^{(\infty)}, \lambda_n^{(\infty)}\}$
constructed along a common diagonal subsequence
(Lemma~\ref{lem:diagonal}).
Without assuming completeness, we establish:
\begin{enumerate}
\item \textbf{Spectral inclusion}
      (Thm.~\ref{thm:spectral-inclusion}, unconditional):
      $\overline{\{\lambda_n^{(\infty)}\}}
      \subset \sigma(\Hlim)\subset[0,1]$.
\item \textbf{Semiclassical spectral transfer}
      (Thm.~\ref{thm:weyl}, under Hyp.~\ref{hyp:SOH}):
      $\lambda_n^{(\infty)}\approx\lambda$ iff
      $n\approx\frac{2}{\pi}c(\lambda)$, where
      $c(\lambda)$ is the finite-$c$ transition
      bandwidth for level $\lambda$.
\item \textbf{Density criterion}
      (Thm.~\ref{thm:density-criterion}):
      density of $\{\lambda_n^{(\infty)}\}$ implies
      $\sigma(\Hlim)=[0,1]$.
\end{enumerate}
Riesz-basis stability and the Bridge theorem
are deferred to Paper~XII.
No unconditional connection to $\zeta(s)$-zeros.
\end{abstract}

\section*{Guide for the Referee}
Two inherited layers:
\begin{itemize}
\item \textbf{Operator layer} (Paper~X):
      $\Hlim$ bounded, self-adjoint, SOT-limit.
\item \textbf{Eigenstructure layer}
      (Paper~IX + Lemma~\ref{lem:diagonal}):
      orthonormal system
      $\{\Phi_n^{(\infty)}, \lambda_n^{(\infty)}\}$
      along a common diagonal subsequence.
\end{itemize}
Paper~XI adds the spectral geometry layer:
spectral inclusion (unconditional),
semiclas\-sical transfer (conditional on Hyp.~\ref{hyp:SOH}),
density reduction.
No completeness assumed; frame-stability
deferred to Paper~XII.

\tableofcontents

%%═══════════════════════════════════════════════════════════════════
\section{Main Results}
\label{sec:main}
%%═══════════════════════════════════════════════════════════════════

\begin{theorem}[Main results of Paper~XI]
\label{thm:main}
Let $\Hlim$ and $\{\Phi_n^{(\infty)},\lambda_n^{(\infty)}\}$
be as in Lemma~\ref{lem:diagonal}.
\begin{enumerate}[\rm(A)]
\item \textbf{(Eigenvalue equation,
      Lem.~\ref{lem:eigenvalue-eq}.)}
      $\Hlim\,\Phi_n^{(\infty)}
      =\lambda_n^{(\infty)}\,\Phi_n^{(\infty)}$
      for all $n\ge 0$.
      Uses only SOT, eigenvalue rate, orthonormality;
      no completeness, no resolvent, no spectral theorem.
\item \textbf{(Spectral inclusion,
      Thm.~\ref{thm:spectral-inclusion}.)}
      $\overline{\{\lambda_n^{(\infty)}\}}
      \subset\sigma(\Hlim)\subset[0,1]$.
      Unconditional.
\item \textbf{(Semiclassical transfer,
      Thm.~\ref{thm:weyl}.)}
      Under Hyp.~\ref{hyp:SOH}:
      $\lambda_n^{(\infty)}\approx\lambda$
      iff $n\approx\frac{2}{\pi}c(\lambda)$,
      where $c(\lambda)$ is the
      finite-$c$ transition bandwidth
      (Def.~\ref{def:clambda}).
\item \textbf{(Density criterion,
      Thm.~\ref{thm:density-criterion}.)}
      Density of $\{\lambda_n^{(\infty)}\}$ in $[0,1]$
      implies $\sigma(\Hlim)=[0,1]$.
\end{enumerate}
(A)--(B) unconditional;
(C) conditional on Hyp.~\ref{hyp:SOH};
(D) conditional on Conj.~\ref{conj:density}.
\end{theorem}

%%═══════════════════════════════════════════════════════════════════
\section{Setup and Inherited Results}
\label{sec:setup}
%%═══════════════════════════════════════════════════════════════════

\subsection{Common diagonal subsequence}

\begin{lemma}[Common diagonal subsequence]
\label{lem:diagonal}
There exists $c_k\nearrow\infty$ such that
for every $n\ge 0$ simultaneously:
\begin{enumerate}[\rm(i)]
\item $\Phi_n^{(c_k)}\rightharpoonup\Phi_n^{(\infty)}$
      weakly in $L^2(\mathbb{R})$;
\item $\lambda_n^{(c_k)}\to\lambda_n^{(\infty)}$;
\item $\mathcal{H}_{c_k}
      \xrightarrow{\mathrm{SOT}}\Hlim$.
\end{enumerate}
$\{\Phi_n^{(\infty)}\}_{n\ge 0}$ is orthonormal.
\end{lemma}

\begin{proof}
\textbf{Step~1.}
$\|\Phi_n^{(c)}\|=1$; Banach--Alaoglu.

\textbf{Step~2 (Cantor diagonal).}
Extract $c_k^{(n)}$ inductively;
set $c_k:=c_k^{(k)}$.
For fixed $n$, $(c_k)_{k\ge n}$ is a subsequence
of $c_k^{(n)}$, giving weak convergence (i).
Eigenvalue convergence (ii) from
$|\lambda_n^{(c)}-\lambda_n^{(\infty)}|\le C_\kappa c^{-1/4}$
(Paper~VIII, Cor.~4.3).

\textbf{Step~3.}
(iii) holds along any sequence (Paper~X, Lem.~2.2).

\textbf{Step~4 (Orthonormality).}
$\ip{\Phi_m^{(c_k)}}{\Phi_n^{(c_k)}}=0$ for $m\ne n$;
weak convergence gives
$\ip{\Phi_m^{(\infty)}}{\Phi_n^{(\infty)}}=0$.
$\|\Phi_n^{(\infty)}\|=1$: Paper~IX, Prop.~2.2.
\end{proof}

\begin{remark}[Scope]
All results along the fixed $c_k$.
Uniqueness up to phase: Problem~\ref{prob:uniqueness}.
\end{remark}

\subsection{Inherited results}

\begin{enumerate}[(XI.1)]
\item\label{XI1}
  \textbf{Eigenvalue rate} (Paper~VIII, Cor.~4.3):
  $|\lambda_n^{(c)}-\lambda_n^{(\infty)}|
  \le C_\kappa c^{-1/4}$, $n\le\kappa c$.
\item\label{XI2}
  \textbf{Orthonormal system}
  (Lem.~\ref{lem:diagonal} + Paper~IX, Prop.~2.2):
  $\Phi_n^{(c_k)}\rightharpoonup\Phi_n^{(\infty)}$;
  $\|\Phi_n^{(\infty)}\|=1$;
  $\ip{\Phi_m^{(\infty)}}{\Phi_n^{(\infty)}}=\delta_{mn}$.
\item\label{XI3}
  \textbf{SOT-convergence} (Paper~X, Lem.~2.2):
  $\mathcal{H}_{c_k}
  \xrightarrow{\mathrm{SOT}}\Hlim$
  (SOT holds pointwise on all of $H$;
boundedness $\|\mathcal{H}_{c_k}\|\le 1$).
\item\label{XI4}
  \textbf{Non-cancellation} (Paper~VIII, Thm.~3.3):
  $|\Phihat_n^{(c)}|\ge c_1 c^{-1/2}$ on
  $E_c$ with $|E_c|\ge c_*c^{-1/2}$.
\item\label{XI5}
  \textbf{Form rate} (Paper~VIII, Cor.~4.3):
  $0\le q_{\lim}-q_c\le C_\kappa c^{-1/4}\|{\cdot}\|^2$.
\item\label{XI6}
  \textbf{Slepian counting}
  \cite{Slepian1961,Osipov2013}:
  transition at $n=2c/\pi$;
  $O(\log c)$ eigenvalues in
  $[\varepsilon,1{-}\varepsilon]$ for each $c$.
\end{enumerate}

\begin{remark}[Dependency architecture]
\textbf{Paper~IX:}
orthonormality, norm control, spectral labeling.
\textbf{Paper~X:}
SOT-existence, self-adjointness,
$\sigma(\Hlim)\subset[0,1]$.
\textbf{Paper~XI:}
eigenvalue equation, spectral inclusion,
semiclas\-sical transfer, density reduction.
\end{remark}

\begin{remark}[Why Paper~IX is necessary]
SOT + non-cancellation give $\Phi_n^{(\infty)}\ne 0$
but no global norm control
($E_c\to\emptyset$, no $L^2$-lower bound survives).
$\|\Phi_n^{(\infty)}\|=1$ is an output of Paper~IX.
\end{remark}

%%═══════════════════════════════════════════════════════════════════
\section{Spectral Inclusion}
\label{sec:spectral-inclusion}
%%═══════════════════════════════════════════════════════════════════

\begin{lemma}[Non-vanishing]
\label{lem:nonvanish-limit}
$\|\Phi_n^{(\infty)}\|=1$.
\end{lemma}
\begin{proof}Orthonormality \eqref{XI2}.\end{proof}

\begin{lemma}[Strong convergence]
\label{lem:strong-conv}
Let $x_k\rightharpoonup x$ weakly in a Hilbert
space with $\|x_k\|=1=\|x\|$.
Then $x_k\to x$ strongly.
\end{lemma}
\begin{proof}
$\|x_k-x\|^2
=2-2\,\mathrm{Re}\ip{x_k}{x}\to 0$
by $\ip{x_k}{x}\to\ip{x}{x}=1$.
\end{proof}

\begin{remark}[Independence of norm sources]
$\|\Phi_n^{(c_k)}\|=1$ (unit eigenvectors);
$\|\Phi_n^{(\infty)}\|=1$ (Paper~IX orthonormality).
The two inputs are independent;
norm equality is not derived from weak convergence.
\end{remark}

\begin{lemma}[Eigenvalue equation]
\label{lem:eigenvalue-eq}
$\Hlim\,\Phi_n^{(\infty)}
=\lambda_n^{(\infty)}\,\Phi_n^{(\infty)}$.
\end{lemma}
\begin{proof}
\textbf{Step~1.}
SOT \eqref{XI3} applied pointwise to
$\Phi_n^{(\infty)}\in L^2$:
$\mathcal{H}_{c_k}\,\Phi_n^{(\infty)}
\to\Hlim\,\Phi_n^{(\infty)}$.

\textbf{Step~2.}
$\mathcal{H}_{c_k}\,\Phi_n^{(\infty)}
=\mathcal{H}_{c_k}(\Phi_n^{(\infty)}-\Phi_n^{(c_k)})
+\lambda_n^{(c_k)}\,\Phi_n^{(c_k)}$.

\textbf{Step~3.}
Lemma~\ref{lem:strong-conv} + \eqref{XI2}:
$\Phi_n^{(c_k)}\to\Phi_n^{(\infty)}$ strongly.

\textbf{Step~4.}
$\|\mathcal{H}_{c_k}(\Phi_n^{(\infty)}-\Phi_n^{(c_k)})\|
\le\|\Phi_n^{(\infty)}-\Phi_n^{(c_k)}\|\to 0$
($\|\mathcal{H}_{c_k}\|\le 1$).

\textbf{Step~5.}
$\|\lambda_n^{(c_k)}\Phi_n^{(c_k)}
-\lambda_n^{(\infty)}\Phi_n^{(\infty)}\|
\to 0$ by \eqref{XI1} and Step~3.
\end{proof}

\begin{remark}[Why SOT suffices here]
In general SOT-limits do not preserve spectra.
The argument works because:
(i) eigenvectors are unit vectors;
(ii) limit norm $=1$ is independent (Paper~IX);
(iii) eigenvalue rate is quantitative ($c^{-1/4}$).
No resolvent or compactness needed.
\end{remark}

\begin{theorem}[Spectral inclusion]
\label{thm:spectral-inclusion}
\begin{enumerate}[\rm(a)]
\item $\lambda_n^{(\infty)}\in\sigma_p(\Hlim)$
      for each $n\ge 0$.
\item $\overline{\{\lambda_n^{(\infty)}\}}
      \subset\sigma(\Hlim)$.
\end{enumerate}
\end{theorem}
\begin{proof}
(a) Lemma~\ref{lem:eigenvalue-eq};
$\Phi_n^{(\infty)}\ne 0$.
(b) $\sigma(\Hlim)$ closed.
\end{proof}

\begin{remark}[Point, approximate, and full spectrum]
\label{rem:sigma-app}
The three spectral inclusions for $\Hlim$ are:
\[
  \overline{\{\lambda_n^{(\infty)}\}}
  \;\subset\; \sigapp(\Hlim)
  \;\subset\; \sigma(\Hlim)
  \;\subset\; [0,1].
\]
The outer bounds are unconditional (this paper and
Paper~X resp.).
The middle term $\sigapp(\Hlim)$ contains the
approximate point spectrum:
$\lambda\in\sigapp(\Hlim)$ iff there exist
unit vectors $f_n$ with
$(\Hlim-\lambda)f_n\to 0$.
Since each $\lambda_n^{(\infty)}$ is a genuine
eigenvalue (part~(a)), we have
$\sigma_p(\Hlim)\subset\sigapp(\Hlim)$,
hence
$\overline{\{\lambda_n^{(\infty)}\}}
\subset\sigapp(\Hlim)$.
The equality $\sigapp(\Hlim)=\sigma(\Hlim)$
would follow from self-adjointness of $\Hlim$
on a core domain, i.e.\ essential self-adjointness
(Paper~IX, conditional on completeness).
Paper~XI's unconditional contribution is
$\sigma_p(\Hlim)$ non-empty and explicitly identified,
without any completeness or domain assumption.
\end{remark}

\begin{remark}[Bounds on $\sigma(\Hlim)$]
Lower: $\overline{\{\lambda_n^{(\infty)}\}}
\subset\sigma(\Hlim)$ (above).
Upper: $\sigma(\Hlim)\subset[0,1]$ (Paper~X).
Equality: Section~\ref{sec:density}.
\end{remark}

%%═══════════════════════════════════════════════════════════════════
\section{Semiclassical Spectral Transfer}
\label{sec:weyl}
%%═══════════════════════════════════════════════════════════════════

\begin{definition}[Finite-$c$ transition bandwidth]
\label{def:clambda}
For $\lambda\in(0,1)$, the
\emph{transition bandwidth} $c(\lambda)$ is
defined by the finite-$c$ Slepian data:
\begin{equation}\label{eq:clambda}
  c(\lambda)
  := \inf\bigl\{c > 0 :
     \lambda_{\lfloor 2c/\pi\rfloor}^{(c)}
     \le \lambda\bigr\}.
\end{equation}
Intuitively: $c(\lambda)$ is the bandwidth at which
the Slepian transition eigenvalue crosses level
$\lambda$.
This is defined entirely from the finite-$c$
family, independently of the limit objects
$\lambda_n^{(\infty)}$.
\end{definition}

\begin{remark}[Well-definedness of $c(\lambda)$]
\label{rem:clambda-welldefined}
For each fixed $c>0$, the Slepian counting law
\eqref{XI6} guarantees that
$\lambda_{\lfloor 2c/\pi\rfloor}^{(c)}$
is in the transition window;
as $c\to\infty$ the transition eigenvalue
sweeps $(0,1)$ by IVT.
Hence the infimum in \eqref{eq:clambda} is
finite for all $\lambda\in(0,1)$.
Under Hypothesis~\ref{hyp:SOH},
$c(\lambda)$ is strictly increasing in $\lambda$.
\end{remark}

\begin{definition}[Counting function]
\label{def:counting}
$N(\lambda):=\#\{n\ge 0:\lambda_n^{(\infty)}\le\lambda\}$
(non-decreasing, right-continuous;
finite for all $\lambda\in[0,1)$
under Hyp.~\ref{hyp:SOH}).
\end{definition}

The Slepian--Pollak--Widom theorem
\cite{Slepian1961,LandauPollak1962,Osipov2013,Widom1964}
gives, for each $\varepsilon>0$:
\begin{equation}\label{eq:slepian-count}
  \#\{n:\lambda_n^{(c)}\in[\varepsilon,1{-}\varepsilon]\}
  =O(\log c),
\end{equation}
transition window centred at $n=2c/\pi$.

\begin{hypothesis}[Spectral Ordering (SOH)]
\label{hyp:SOH}
$m\ne n\implies\lambda_m^{(\infty)}\ne\lambda_n^{(\infty)}$.
The limiting eigenvalue sequence is injective
(no spectral collisions).
\end{hypothesis}

\begin{remark}[SOH and gap control]
\label{rem:SOH-status}
SOH would follow from the uniform gap bound:
\begin{equation}\label{eq:gap-lower}
  |\lambda_n^{(c)}-\lambda_{n+1}^{(c)}|
  \ge C_\kappa c^{-3/4},\quad n\le\kappa c,
\end{equation}
combined with rate \eqref{XI1}.
Estimate \eqref{eq:gap-lower} is the lower bound
of Problem~\ref{prob:rate-sharp};
SOH (Problem~\ref{prob:SOH}) is downstream.
\end{remark}

\begin{theorem}[Semiclassical spectral transfer]
\label{thm:weyl}
Under Hypothesis~\ref{hyp:SOH},
the limiting eigenvalue $\lambda_n^{(\infty)}$
and the finite-$c$ transition bandwidth
$c(\lambda)$ (Def.~\ref{def:clambda})
are related by a genuine asymptotic equivalence:
\begin{enumerate}[\rm(a)]
\item \textbf{(Transfer law.)} For $\lambda\in(0,1)$:
\begin{equation}\label{eq:transfer}
  \lambda_n^{(\infty)}\approx\lambda
  \;\iff\;
  n \approx \tfrac{2}{\pi}\,c(\lambda)
  \quad\text{as }\lambda\to 1^-.
\end{equation}
Precisely:
$N(\lambda)
:= \#\{n:\lambda_n^{(\infty)}\le\lambda\}
\sim \frac{2}{\pi}c(\lambda)$.
\item \textbf{(Transition window.)} For
      $\lambda\in[\varepsilon,1{-}\varepsilon]$:
      $N(\lambda)-N(\lambda-\varepsilon)
      =O(\log c(\lambda))$.
\item \textbf{(Spectral floor.)}
      $N(\lambda)=0$ for
      $\lambda<\lambda_0^{(\infty)}$.
\end{enumerate}
\end{theorem}

\begin{proof}
(a) At bandwidth $c(\lambda)$,
the Slepian transition sits at index
$n_0:=\lfloor 2c(\lambda)/\pi\rfloor$.
By \eqref{XI6}, $\lambda_{n_0}^{(c(\lambda))}\approx\lambda$.
By \eqref{XI1},
$|\lambda_{n_0}^{(c(\lambda))}-\lambda_{n_0}^{(\infty)}|
\le C_\kappa c(\lambda)^{-1/4}\to 0$.
Hence $\lambda_{n_0}^{(\infty)}\approx\lambda$,
with $n_0\sim 2c(\lambda)/\pi$.
This is the genuine content of \eqref{eq:transfer}:
the Slepian bandwidth index transfers to the
limiting spectral scale via the rate \eqref{XI1}.
(b) From \eqref{eq:slepian-count} and \eqref{XI1}.
(c) No $\lambda_n^{(\infty)}$ below
$\lambda_0^{(\infty)}=\lim_k\lambda_0^{(c_k)}$.
\end{proof}

\begin{remark}[Non-circularity]
\label{rem:noncircular}
$c(\lambda)$ is defined from finite-$c$ data
(Def.~\ref{def:clambda}), independently of
$\{\lambda_n^{(\infty)}\}$.
The theorem is non-trivial: it asserts that the
limit eigenvalue index $N(\lambda)$ coincides
asymptotically with the Slepian bandwidth
index $2c(\lambda)/\pi$, a transfer across two
apriori unrelated scales.
\end{remark}

\begin{remark}[Not a classical Weyl law]
$\Hlim$ is not compact; $N(\lambda)$ is finite
only under SOH.
This is a \emph{semiclassical index transfer},
not a counting law for a compact operator.
\end{remark}

%%═══════════════════════════════════════════════════════════════════
\section{Density Criterion}
\label{sec:density}
%%═══════════════════════════════════════════════════════════════════

\begin{theorem}[Density criterion]
\label{thm:density-criterion}
If $\{\lambda_n^{(\infty)}\}$ is dense in $[0,1]$,
then $\sigma(\Hlim)=[0,1]$.
\end{theorem}
\begin{proof}
Each $\lambda_n^{(\infty)}\in\sigma(\Hlim)$
(Thm.~\ref{thm:spectral-inclusion}(a));
$\sigma(\Hlim)$ closed; density gives
$[0,1]\subset\sigma(\Hlim)$.
Reverse: Paper~X.
\end{proof}

\begin{corollary}
Under density, $\sigma(\Hlim)$ has no isolated
points in $(0,1)$.
\end{corollary}

%%═══════════════════════════════════════════════════════════════════
\section{The Density Question}
\label{sec:density-question}
%%═══════════════════════════════════════════════════════════════════

\begin{conjecture}[Density]
\label{conj:density}
$\{\lambda_n^{(\infty)}\}$ is dense in $[0,1]$.
\end{conjecture}

\subsection{Evidence and reformulation}

By \eqref{eq:slepian-count}, for each $c$
there are $O(\log c)$ eigenvalues in
$[\varepsilon,1{-}\varepsilon]$.
By \eqref{XI1}, each lies within $C_\kappa c^{-1/4}$
of $\lambda_n^{(\infty)}$.

\begin{remark}[Two-parameter reformulation]
\label{rem:density-gap}
Density of $\{\lambda_n^{(\infty)}\}$ is equivalent to:
\begin{quote}
  For every open $I\subset(0,1)$,
  there exist $k_j\to\infty$ and $n_j\ge 0$ with
  $\lambda_{n_j}^{(c_{k_j})}\in I$ for all $j$.
\end{quote}
The IVT guarantees that every $\lambda\in(0,1)$
is \emph{attained} by some $\lambda_n^{(c)}$
for each fixed $c$.
The difficulty is not the existence of
transition eigenvalues in $I$,
but their \emph{uniform distribution} across the
transition band as $c$ varies:
the $O(\log c)$ window must sweep every
subinterval of $(0,1)$ densely as $c\to\infty$.
This is a \emph{moving-window density problem},
analogous to equidistribution questions in
semiclas\-sical analysis;
it is not currently available from Papers~VI--VIII.
\end{remark}

\begin{proposition}[Non-emptiness near $1$]
$\{\lambda_n^{(\infty)}\}\cap[1{-}\delta,1)
\ne\emptyset$ for all $\delta\in(0,1)$.
\end{proposition}
\begin{proof}\eqref{XI6} + \eqref{XI1}.\end{proof}

Numerical tables \cite{Slepian1961,Osipov2013,
LandauPollak1962} show transition eigenvalues
filling $[0,1]$ densely for finite $c$.

%%═══════════════════════════════════════════════════════════════════
\section{Open Problems}
\label{sec:open-problems}
%%═══════════════════════════════════════════════════════════════════

\begin{problem}[Density]
\label{prob:density}
Prove Conj.~\ref{conj:density};
equivalent to the two-parameter coverage
of Rem.~\ref{rem:density-gap}.
\end{problem}

\begin{problem}[Spectral Ordering]
\label{prob:SOH}
Prove Hyp.~\ref{hyp:SOH};
follows from \eqref{eq:gap-lower}
(Rem.~\ref{rem:SOH-status}).
\end{problem}

\begin{problem}[Uniqueness]
\label{prob:uniqueness}
$\Phi_n^{(\infty)}$ independent of the diagonal
subsequence, up to phase.
\end{problem}

\begin{problem}[Form rate sharpness]
\label{prob:rate-sharp}
$\sup_f(q_{\lim}-q_c)(f,f)/\|f\|^2
\asymp c^{-1/4}$?
Lower bound requires \eqref{eq:gap-lower}.
\end{problem}

\begin{problem}[Completeness]
\label{prob:completeness}
$\{\Phi_n^{(\infty)}\}$ complete in $L^2$?
Equivalent to $\sigapp(\Hlim)=\sigma(\Hlim)$.
\end{problem}

\begin{problem}[Riesz-basis stability]
\label{prob:riesz}
$\{\Phi_n^{(\infty)}\}$ Riesz basis for its span?
Implies completeness, $\HSOT=\overline{\Hspectral}$.
Primary target: Paper~XII.
\end{problem}

\begin{remark}[Hierarchy]
\ref{prob:density} + \ref{prob:completeness}:
independent; both needed for Hilbert--P\'olya.
\ref{prob:riesz} resolves completeness + Bridge.
\ref{prob:SOH} $\Leftarrow$ \ref{prob:rate-sharp}.
\ref{prob:uniqueness}: removes subsequential
qualifiers.
\end{remark}

%%═══════════════════════════════════════════════════════════════════
\section{Outlook}
\label{sec:outlook}
%%═══════════════════════════════════════════════════════════════════

Papers~I--XI \emph{reduce} the Hilbert--P\'olya
program to the following independent structural
conditions; no RH implication is claimed without
all conditions being established:

\begin{enumerate}
\item \textbf{Density}
      (Prob.~\ref{prob:density}):
      $\sigma(\Hlim)=[0,1]$.
\item \textbf{Completeness}
      (Prob.~\ref{prob:completeness}):
      $\sigapp(\Hlim)=\sigma(\Hlim)$.
\item \textbf{Spectral ordering}
      (Hyp.~\ref{hyp:SOH}):
      counting law unconditional.
\item \textbf{Uniqueness}
      (Prob.~\ref{prob:uniqueness}):
      removes subsequential qualifiers.
\end{enumerate}

Under all of (1)--(4) together with
Hypothesis~CCM (\cite{CCM2025})
-- which identifies the spectrum of a related
operator with the imaginary parts of non-trivial
zeros of $\zeta(s)$ --
the Riemann Hypothesis would follow.
This implication chain is not yet closed;
the present paper establishes only the
unconditional foundation (A)--(B)
of Theorem~\ref{thm:main}.

Problem~\ref{prob:riesz} would simultaneously
resolve completeness and the Bridge theorem
$\HSOT=\overline{\Hspectral}$;
it is the primary target of Paper~XII.

\medskip\noindent
This is the precise state of the
Hilbert--P\'olya programme as of Paper~XI.

%%%─── Bibliography
\begin{thebibliography}{99}

\bibitem{PaperVIII}
[Author],
\textit{Non-cancellation control, the sharp
$c^{-1/2}$ lower bound, and spectral connection
to the zeros of $\zeta(s)$},
Paper~VIII, preprint, 2026.

\bibitem{PaperIX}
[Author],
\textit{Domain identification, deficiency indices,
and essential self-adjointness of the PSWF
concentration limit operator},
Paper~IX, preprint, 2026.

\bibitem{PaperX}
[Author],
\textit{Mosco convergence, resolvent stability,
and the bridge theorem for PSWF concentration operators},
Paper~X, preprint, 2026.

\bibitem{Slepian1961}
D.~Slepian and H.~O.~Pollak,
\textit{Prolate spheroidal wave functions,
Fourier analysis and uncertainty~I},
Bell System Tech.\ J.\ \textbf{40} (1961), 43--63.

\bibitem{LandauPollak1962}
H.~J.~Landau and H.~O.~Pollak,
\textit{Prolate spheroidal wave functions,
Fourier analysis and uncertainty~III},
Bell System Tech.\ J.\ \textbf{41} (1962), 1295--1336.

\bibitem{Osipov2013}
A.~Osipov, V.~Rokhlin, and H.~Xiao,
\textit{Prolate Spheroidal Wave Functions of Order Zero},
Springer, New York, 2013.

\bibitem{Widom1964}
H.~Widom,
\textit{Asymptotic behavior of the eigenvalues
of certain integral equations~II},
Arch.\ Ration.\ Mech.\ Anal.\ \textbf{17}
(1964), 215--229.

\bibitem{Kato1966}
T.~Kato,
\textit{Perturbation Theory for Linear Operators},
Springer, Berlin, 1966.

\bibitem{ReedSimon1975}
M.~Reed and B.~Simon,
\textit{Methods of Modern Mathematical Physics,
Vols.~I--II},
Academic Press, New York, 1972--1975.

\bibitem{Mosco1969}
U.~Mosco,
\textit{Convergence of convex sets and of
solutions of variational inequalities},
Adv.\ Math.\ \textbf{3} (1969), 510--585.

\bibitem{KuwaeShioya2003}
K.~Kuwae and T.~Shioya,
\textit{Convergence of spectral structures},
Comm.\ Anal.\ Geom.\ \textbf{11} (2003), 599--673.

\bibitem{CCM2025}
A.~Connes, C.~Consani, and H.~Moscovici,
\textit{Zeta Spectral Triples},
arXiv:2511.22755 [math.NT], preprint, 2025.

\end{thebibliography}

\end{document}
